\section*{Chapter \ref{ch:m3:when-markets-break-together} --- When Markets Break Together}

\begin{solution}{exo:m3:when-markets-break-together:1}
$10\text{ million}/0.125 = $ USD~80 million.
\end{solution}

\begin{solution}{exo:m3:when-markets-break-together:2}
In a \omterm{def:m3:when-markets-break-together:flight}{flight to quality} investors sell risky assets and buy the safest, raising government bond prices; in a dash for cash they sell whatever can be
sold, safe assets included, to raise cash, and government bond prices fall with the rest.
\end{solution}

\begin{solution}{exo:m3:when-markets-break-together:3}
Because it is held by the same leveraged investors who need cash: in a dash for cash assets are sold for their liquidity, not for their outlook, and
the correlations estimated in calm break down.
\end{solution}

\begin{solution}{exo:m3:when-markets-break-together:4}
USD~223.6 billion against 124.4 billion without price impact: about USD~99.3 billion, 44\%, comes from the prices that forced sales move, and the losses
they inflict.
\end{solution}

\begin{solution}{exo:m3:when-markets-break-together:5}
About $9 \times 0.64\% = 5.8\%$ of its value in seven trading days, for the safest asset in the world.
\end{solution}

\begin{solution}{exo:m3:when-markets-break-together:6}
Value-at-risk aggregates positions with a correlation matrix; if correlations move towards one in stress, offsets estimated in calm disappear and the
losses add up, so the calm estimate understates the loss exactly when it matters.
\end{solution}

\begin{solution}{exo:m3:when-markets-break-together:7}
At 4 times leverage the shock creates no shortfall: equity covers the higher margins, nothing is sold, and the spiral never starts. Leverage decides
whether the shock is absorbed or amplified.
\end{solution}

\begin{solution}{exo:m3:when-markets-break-together:8}
Diversification by asset class does not help when the holders, the financing and the margin rules are common: in a joint stress the firm's margins rise
everywhere at once, and assets held by the same leveraged investors are sold together.
\end{solution}

\begin{solution}{pb:m3:when-markets-break-together:1}
\textbf{1.} About USD~25.2 billion. \textbf{2.} About USD~223.6 billion. \textbf{3.} 8.88. \textbf{4.} 4.94, the reciprocal of the average margin rate
after margins rose (\cref{prop:m3:when-markets-break-together:multiplier}): each dollar sold releases only the margin rate of requirement. \textbf{5.} They
fall with it, to about 93 and 80 on day 20, because they are sold to meet margin calls on the whole portfolio. \textbf{6.} 8.73: the result does not depend
on the step, because deleveraging is iterated to a fixed point in each step. \textbf{7.} In this run the shock is absorbed without forced sales: the spiral
needs margins that rise. \textbf{8.} Forced sales stay high but the correlation of the untouched assets with the shocked one stays near zero: common holders
are what joins the markets. \textbf{9.} From 0.30 to 0.99. \textbf{10.} The \omterm{def:m3:when-markets-break-together:spirals}{margin spiral} starts it here and price impact nearly doubles it. \textbf{11.}
Treasuries sold with equities, yields up sharply, market depth collapsing, and the basis between cash bonds and futures dislocated. \textbf{12.} Raised
initial margins by roughly USD~300 billion over the month, with variation margin calls peaking at USD~140 billion on 9 March. \textbf{13.} It contained the
cleared swaps; the crisis ran through uncleared derivatives and short-term funding, and through haircuts on collateral of every kind. \textbf{14.} The
gilt episode: leveraged holders, margin calls rising with the move, forced sales into a falling market, until a buyer of last resort stepped in.
\textbf{15.} Central banks bought the assets being sold and supplied liquidity, which breaks both spirals from outside. \textbf{16.} They must cover the risk
they see, but margins that jump with volatility are procyclical; stable, conservative margins and liquidity buffers built in calm, or anti-procyclicality
floors, reduce the spiral. \textbf{17.} From the margin increase on all its positions at once in a stress calibrated on episodes like March 2020, plus
funding that may not roll, held in cash or assets that stay liquid in a dash for cash. \textbf{18.} Leverage and crowding in its markets, margin changes
announced by clearing houses, funding spreads and repo haircuts, market depth, and correlations between assets that are usually independent.
\textbf{19.} \emph{Named result:} a forced-sale multiplier of 8.88 dollars of forced sales per dollar of initial margin shortfall (8.73 with eight steps a
day), against 4.94 without price impact. \textbf{20.} Because they share holders, financing and margin rules, and in a crisis those links carry the selling
from one market to all.
\end{solution}

\begin{solution}{iq:m3:when-markets-break-together:1}
A \omterm{def:m3:when-markets-break-together:spirals}{loss spiral}: losses force sales, sales lower prices, lower prices cause losses. A \omterm{def:m3:when-markets-break-together:spirals}{margin spiral}: volatility raises margins, margins force sales, sales raise
volatility. Together they turn a shock into forced selling across everything leveraged holders own.
\iqlookfor{both feedbacks and why they cross markets.}
\end{solution}

\begin{solution}{iq:m3:when-markets-break-together:2}
In the dash for cash, holders sold what they could sell to raise dollars; leveraged non-bank investors and foreign holders sold Treasuries; dealers could not
absorb the flow, depth collapsed and yields rose.
\iqlookfor{liquidity demand, not credit risk.}
\end{solution}

\begin{solution}{iq:m3:when-markets-break-together:3}
From past stress windows rather than full samples, with correlations pushed towards one for assets held by the same leveraged investors, and by simulation of
forced selling on known holdings, checked against the episodes of 2008, 2020 and 2022.
\iqlookfor{conditional correlations and structural reasoning.}
\end{solution}

\begin{solution}{iq:m3:when-markets-break-together:4}
Scenario: margins up by the worst historical increase on every venue at once, prices down jointly, funding lines not rolled, \omterm{def:m3:blockchains-for-traders:stable}{stablecoin} \omterm{def:m3:blockchains-for-traders:stable}{depeg} and venue
withdrawals slowed; compute the cash needed day by day against the cash and liquid collateral available, including transfer times between venues; set buffers
and limits so the firm survives it without forced sales.
\iqlookfor{joint margins, transfer times and a day-by-day cash ladder.}
\end{solution}

\begin{solution}{iq:m3:when-markets-break-together:5}
The basis widens against it and futures margins rise at once; at 50 times leverage a small move exhausts its equity and it must sell Treasuries and buy back
futures into a falling market, adding to the dash for cash; its prime broker may raise haircuts or pull financing.
\iqlookfor{margin, haircut and forced unwind.}
\end{solution}

\begin{solution}{iq:m3:when-markets-break-together:6}
Holders with positions and equity, margin rules on measured volatility, a price-impact function, a shock; deleverage to a fixed point within each step and run
with several step sizes, checking that results converge; ablate each mechanism and check that the effects go the right way.
\iqlookfor{fixed points, convergence in the step, and ablations.}
\end{solution}
